%------------------------------------------------------------------------------%
\documentclass[fleqn,utf8,aspectratio=169,12pt]{beamer}

\usepackage{integracao_e_series}

\title{Integração por Partes}

%------------------------------------------------------------------------------%
\begin{document}

\addtitle

%------------------------------------------------------------------------------%
\section{Integração por partes}

%------------------------------------------------------------------------------%
\begin{proofframe}{Justificativa}
  \onslide<+->{Assumindo que
  \[
    F'(x) = f(x)
    \qquad\qquad
    G'(x) = g(x)
  \]}
  \onslide<+->{temos}
  \begin{align*}
    \onslide<2->{\big(F(x)G(x)\big)'
      & \;=\; f(x)G(x)
        \;+\; {\color<3>{blue}F(x)g(x)} \\[3mm]}
    \onslide<+->{{\color<3>{blue}F(x)g(x)}
      & \;=\; \big(F(x)G(x)\big)'
        \;-\; f(x)G(x) \\[2mm]}
    \onslide<+->{\int F(x)g(x)\,dx
      & \;=\; {\color<5>{blue}\int \big(}F(x)G(x){\color<5>{blue}\big)'\,dx}
        \;-\; \int f(x)G(x)\,dx \\[2mm]}
    \onslide<+->{\int F(x) \, g(x)\,dx
      & \;=\;         F(x) \, G(x)
        \;-\; \int    G(x) \, f(x)\,dx }
  \end{align*}
\end{proofframe}

%------------------------------------------------------------------------------%
\begin{frame}{Integração por Partes}
  Integração por Partes
  \[
    \int       {\color<4->{blue}F(x)}   \, {\color<2->{magenta}g(x)}\,dx
    \;=\;      {\color<4->{blue}F(x)}   \, {\color<2->{magenta}G(x)}
    \;-\; \int {\color<2->{magenta}G(x)}\, {\color<4->{blue}f(x)}\,dx
  \]

  \pause
  Integramos
  \tabto{21mm} {\color<2->{magenta}\(g(x)\)}
  \pause
  \tabto{33mm} -- \;
  essa função não deve ``complicar'' ao ser integrada

  \pause
  Derivamos
  \tabto{21mm} {\color<4->{blue}\(F(x)\)}
  \pause
  \tabto{33mm} -- \;
  essa função deve ``simplificar'' ao ser derivada
\end{frame}

%------------------------------------------------------------------------------%
\begin{frame}{Integração por Partes}
  Mnemônico
  {\large
  \[
    \int {\color{blue}u}\, {\color{magenta}dv}
    \;=\;
    {\color{blue}u}\,{\color{magenta}v}
    \;-\;
    \int {\color{magenta}v}\,{\color{blue}du}
  \]}
  \begin{align*}
    {\color{blue}u}  & = F(x)     & {\color{magenta}dv} & = g(x) dx & \hspace{1cm}~ \\
    {\color{blue}du} & = f(x)\,dx & {\color{magenta}v}  & = \int g(x) dx = G(x)
  \end{align*}
\end{frame}

%------------------------------------------------------------------------------%
\begin{frame}{Integração por Partes para integrais definidas}
  Integrais definidas
  {\large
  \[
    {\color{blue}\int_a^b} u\, dv
    \;=\; u\,v\, {\color{blue}\evalat{a}{b}}
    \;-\; {\color{blue}\int_a^b} v\,du
  \]}
\end{frame}

%------------------------------------------------------------------------------%
\section{Exemplos}

%------------------------------------------------------------------------------%
\begin{question}
  Calcule a integral
  \(
    \int x \sin(x)\,dx
  \)
\end{question}

%------------------------------------------------------------------------------%
\begin{resolution}{Solução}
  \onslide<+->{Queremos calcular \;
  \(
    H(x) = \int x \sin(x)\,dx
  \)}

  \onslide<+->{usamos \;
  \(
    \int {\color{blue}u}{\color{magenta}dv}
    \;=\;
    {\color{blue}u}{\color{magenta}v}
    \;-\;
    \int {\color{magenta}v} {\color{blue}du}
  \)}
  \begin{align*}
    \onslide<+->{{\color{blue}u}  & =  x & {\color{magenta}dv} & = \sin(x) dx & \hspace{1cm}~ \\}
    \onslide<+->{{\color{blue}du} & = dx & {\color{magenta}v}  & = \int \sin(x) dx = -\cos(x)}
  \end{align*}
  \onslide<+->{temos
  \[
    H(x)
    \;=\; {\color{blue}x} \, \big({\color{magenta}-\cos(x)}\big)
    \;-\; \int \big({\color{magenta}-\cos(x)}\big) {\color{blue}dx}
  \]}
\end{resolution}

%------------------------------------------------------------------------------%
\begin{resolution}{Solução}
  \begin{align*}
    \onslide<+->{H(x) & = \int x\sin(x)dx                                \\}
    \onslide<+->{     & = x\big(-\cos(x)\big)- \int \big(-\cos(x)\big)dx \\}
    \onslide<+->{     & = -x \cos(x) + \int \cos(x) dx                   \\[1mm]}
    \onslide<+->{     & = -x \cos(x) + \sin(x) + c}
  \end{align*}
\end{resolution}

%------------------------------------------------------------------------------%

%------------------------------------------------------------------------------%
\begin{question}
  Calcule a integral
  \(
    \int t \ln t \, dt
  \)
\end{question}

%------------------------------------------------------------------------------%
\begin{resolution}{Solução}
  \onslide<+->{Queremos calcular \;
  \(
    R(t) = \int t \ln t \, dt
  \)}

  \onslide<+->{usamos \;
  \(
    \int u dv \;=\;u v \;-\; \int v du
  \)}
  \begin{align*}
    \onslide<+->{u  & = \ln t           & dv & = t dt & \hspace{3cm}~ \\}
    \onslide<+->{du & = \frac{1}{t}\,dt & v  & = \frac{t^2}{2}}
  \end{align*}
  \onslide<+->{temos \;
  \(
    R(t)
    \;=\; \ln(t) \frac{t^2}{2}
    \;-\; \int \frac{t^2}{2} \frac{1}{t}\,dt
  \)}
\end{resolution}

%------------------------------------------------------------------------------%
\begin{resolution}{Solução}
  \begin{align*}
    \onslide<+->{R(t)   = \int t \ln t \, dt                                        \;}
    \onslide<+->{     & = \ln(t) \frac{t^2}{2} - \int \frac{t^2}{2} \frac{1}{t}\,dt \\}
    \onslide<+->{     & =  \frac{t^2}{2} \ln(t) - \frac{1}{2} \int t dt             \\}
    \onslide<+->{     & =  \frac{t^2}{2} \ln(t) - \frac{1}{2} \frac{t^2}{2} + k     \\}
    \onslide<+->{     & =  \frac{t^2}{2} \left(\ln(t) - \frac{1}{2}\right) + k      \\}
    \onslide<+->{     & =  \frac{t^2}{2} \ln\left(\frac{t}{\sqrt{e}}\right) + k                              }
  \end{align*}
\end{resolution}

%------------------------------------------------------------------------------%

%------------------------------------------------------------------------------%
\begin{question}
  Calcule a integral
  \(
    \int x e^x \, dx
  \)
\end{question}

%------------------------------------------------------------------------------%
\begin{resolution}{Solução}
  \onslide<+->{
  \[
    F(x) = \int x e^x \, dx
  \]}
  \vspace{-1.5\baselineskip}
  \begin{align*}
    \onslide<+->{u  & = x  & dv & = e^x\,dx & \hspace{5cm}~\\}
    \onslide<+->{du & = dx & v  & = e^x}
  \end{align*}
  \vspace{-1.5\baselineskip}
  \begin{align*}
    \onslide<+->{F(x) & = xe^x - \int e^x dx      }\\[1mm]
    \onslide<+->{     & = xe^x - e^x + c          }\\[1mm]
    \onslide<+->{     & = e^x(x -1) + c           }
  \end{align*}

\end{resolution}

%------------------------------------------------------------------------------%

%------------------------------------------------------------------------------%
\begin{question}
  Calcule
  \(
    \int \ln(x) dx
  \)
  \pause
  \(
    = \int \ln(x) \cdot 1 \; dx
  \)
\end{question}

%------------------------------------------------------------------------------%
\begin{resolution}{Solução}
  \onslide<+->{\[
    F(x) = \int \ln(x) dx
  \]}
  \vspace{-\baselineskip}
  \begin{align*}
    \onslide<+->{u  & = \ln(x)        & dv & = 1 dx  & \hspace{4cm}~ }\\
    \onslide<+->{du & = \frac{1}{x}dx & v  & = x                     }
  \end{align*}
  \begin{align*}
    \onslide<+->{F(x) & = x \ln(x) - \int x\, \frac{1}{x}\,dx}\\
    \onslide<+->{     & = x \ln(x) - \int dx                 }\\
    \onslide<+->{     & = x \ln(x) - x + c                   }
  \end{align*}
\end{resolution}

%------------------------------------------------------------------------------%

%------------------------------------------------------------------------------%
\begin{question}
  Calcule
  \(
    \int t^2 e^t dt
  \)
\end{question}

%------------------------------------------------------------------------------%
\begin{resolution}{Solução}
  \onslide<+->{\[
    F(t) = \int t^2 e^t dt
  \]}
  \vspace{-\baselineskip}
  \begin{align*}
    \onslide<+->{u  & = t^2  & dv & = e^t dt & \hspace{5cm}~ \\}
    \onslide<+->{du & = 2tdt & v  & = e^t}
  \end{align*}
  \onslide<+->{\[
    F(t)  = t^2 e^t - 2 \int t e^t\, dt
  \]}
\end{resolution}

%------------------------------------------------------------------------------%
\begin{resolution}{Solução}
  \onslide<+->{\[
    G(t) = \int t e^t\, dt
  \]}
  \vspace{-\baselineskip}
  \begin{align*}
    \onslide<+->{u  & = t  & dv & = e^t dt & \hspace{5cm}~ \\}
    \onslide<+->{du & = dt & v  & = e^t}
  \end{align*}
  \begin{align*}
    \onslide<+->{G(t) & = t e^t - \int e^t dt  \\}
    \onslide<+->{     & = t e^t - e^t + c}
  \end{align*}
\end{resolution}

%------------------------------------------------------------------------------%
\begin{resolution}{Solução}
  \begin{align*}
    \onslide<+->{F(t) & = t^2 e^t - G(t)                           \\[1mm]}
    \onslide<+->{     & = t^2 e^t - 2\left( te^t - e^t + c \right) \\[1mm]}
    \onslide<+->{     & = t^2 e^t - 2te^t + 2e^t - c               \\[1mm]}
    \onslide<+->{     & = \left(t^2 - 2t + 2\right)e^t + c}
  \end{align*}
\end{resolution}

%------------------------------------------------------------------------------%

%------------------------------------------------------------------------------%
\begin{question}
  Calcule a integral
  \(
    \int e^x \sin(x) dx
  \)
\end{question}

%------------------------------------------------------------------------------%
\begin{resolution}{Solução}
  \onslide<+->{\[
    F(x) = \int e^x \sin(x) dx
  \]}
  \vspace{-\baselineskip}
  \begin{align*}
    \onslide<+->{u  & = e^x   & dv & = \sin(x) dx & \hspace{3cm}~ }\\
    \onslide<+->{du & = e^xdx & v  & = -\cos(x) }
  \end{align*}
  \onslide<+->{\[
    F(x) = -e^x \cos(x)+ \int e^x \cos(x)dx
  \]}
\end{resolution}

%------------------------------------------------------------------------------%
\begin{resolution}{Andando em Círculos}
  \onslide<+->{\[
    G(x) = \int e^x \cos(x) dx
  \]}
  \vspace{-\baselineskip}
  \begin{align*}
    \onslide<+->{u  & = e^x   & dv & = \cos(x) dx & \hspace{3cm}~ }\\
    \onslide<+->{du & = e^xdx & v  & = \sin(x) }
  \end{align*}
  \onslide<+->{\[
    G(x) = e^x \sin(x) - \int e^x \sin(x) dx
  \]}
\end{resolution}

%------------------------------------------------------------------------------%
\begin{resolution}{Solução}
  \begin{align*}
    \onslide<+->{
      F(x) = \int e^x \sin(x) dx
      &    = -    e^x \cos(x) + \int e^x \cos(x) dx }\\
    \onslide<+->{
      G(x) = \int e^x \cos(x) dx
      &    = \hpm {\color<3>{blue}e^x \sin(x) - \int e^x \sin(x) dx} }
  \end{align*}
  \begin{align*}
     \onslide<+->{
      F(x) =
      \int e^x \sin(x) dx
      & \;=\; -e^x \cos(x) + {\color<3>{blue}e^x \sin(x) - \int e^x \sin(x) dx}
    }\\
    \onslide<+->{
      2\int e^x \sin(x) dx
      & \;=\; -e^x \cos(x)+ e^x \sin(x)
    }\\
    \onslide<+->{
      \int e^x \sin(x) dx
      & \;=\; \frac{e^x}{2}\big(\sin(x) -\cos(x)\big) + c
    }
  \end{align*}
\end{resolution}

%------------------------------------------------------------------------------%

%------------------------------------------------------------------------------%
\begin{question}
  Calcule a integral
  \(
    \int_0^1 \frac{y}{\,e^{2y}\,} dy
  \)
\end{question}
%------------------------------------------------------------------------------%
\begin{resolution}{Solução}
  \onslide<+->{Primeiro precisamos encontrar a primitiva}
  \[
    \onslide<+->{F(x) \;=\; \int \frac{y}{\,e^{2y}\,} dy}
    \onslide<+->{     \;=\; \int y \,e^{-2y}\, dy}
  \]
  \vspace{-\baselineskip}
  \begin{align*}
    \onslide<+->{u  & = y  & dv & = e^{-2y} dy & \hspace{3cm}~ }\\
    \onslide<+->{du & = dy &}
    \onslide<+->{v & = \int e^{-2y} dy}
    \onslide<+->{    = -\frac{e^{-2y}}{2}}
  \end{align*}
  \onslide<+->{\[
    F(x)
    \;=\; -\frac{ye^{-2y}}{2}- \int -\frac{e^{-2y}}{2} dy
  \]}
\end{resolution}

%------------------------------------------------------------------------------%
\begin{resolution}{Solução}
\begin{align*}
  \onslide<+->{F(x) & = -\frac{y}{2}e^{-2y} - \int -\frac{e^{-2y}}{2} dy         }\\[1mm]
  \onslide<+->{     & = -\frac{y}{2}e^{-2y} + \frac{1}{2} \int e^{-2y} dy        }\\[1mm]
  \onslide<+->{     & = -\frac{y}{2}e^{-2y} + \frac{1}{2} \frac{e^{-2y}}{-2} + c }\\[1mm]
  \onslide<+->{     & = -\frac{y}{2}e^{-2y} - \frac{1}{4} e^{-2y} + c            }
\end{align*}
\end{resolution}

%------------------------------------------------------------------------------%
\begin{resolution}{Solução}
  \onslide<+->{Usando o Teorema Fundamental do Cálculo}
  \begin{align*}
    \onslide<+->{\int_0^1 \frac{y}{\,e^{2y}\,} dy}
    \onslide<+->{& = F(x) \evalat{0}{1}                                                                }\\[1mm]
    \onslide<+->{& = \left( -\frac{y}{2}e^{-2y} - \frac{1}{4} e^{-2y} \right) \evalat{0}{1}            }\\[1mm]
    \onslide<+->{& = \left( -\frac{1}{2}e^{-2}- \frac{1}{4}e^{-2}\right)- \left(0- \frac{1}{4} \right) }\\[1mm]
    \onslide<+->{& = -\frac{3}{4}e^{-2} + \frac{1}{4} }
    \onslide<+->{\;=\; \frac{1-3e^{-2}}{4}            }
  \end{align*}
\end{resolution}

%------------------------------------------------------------------------------%

%------------------------------------------------------------------------------%
\begin{question}
  Considere a região do primeiro quadrante delimitada por
  \[
    f(x) = x \sin(x)
    \qquad\qquad
    0 \leq x \leq \pi
  \]

  \bigskip
  Encontre a área dessa região
\end{question}

\framedgraphic{Gráfico da função $f(x)=x \sin(x)$}{ex_area_x_sinx}{}

%------------------------------------------------------------------------------%
\begin{resolution}{Cálculo da Área}
  Área
  \[
    A
    = \int_a^b f(x) - g(x) \,dx
    \pause
    = \int_0^{\pi} x\sin(x)\,dx
  \]
  \pause
  Precisamos da primitiva
  \[
    F(x)
    \;=\; \int x\sin(x)dx
    \pause
    \;=\; -x \cos(x) + \sin(x) + c
  \]
\end{resolution}

%------------------------------------------------------------------------------%
\begin{resolution}{Cálculo da Área}
  \begin{align*}
    \onslide<+->{A & = F(x)                             \evalat{0}{\pi} \\[1mm]}
    \onslide<+->{  & = \big( -x \cos(x) + \sin(x) \big) \evalat{0}{\pi} \\[1mm]}
    \onslide<+->{  & = \big(-\pi\cos(\pi) + \sin(\pi)\big)
                     - \big(-  0\cos(0)   + \sin(0)\big)                \\[1mm]}
    \onslide<+->{  & = \big(-\pi(-1) + 0\big)
                     - \big(-  0   + 0\big)                             \\[1mm]}
    \onslide<+->{  & = \pi}
  \end{align*}
\end{resolution}

%------------------------------------------------------------------------------%

%------------------------------------------------------------------------------%
\begin{question}
  Considere a região do primeiro quadrante delimitada por
  \[
    f(x) = x \sin(x)
    \qquad\qquad
    0 \leq x \leq \pi
  \]

  \bigskip
  Encontre o volume do sólido de revolução gerado pela \\
  rotação dessa região em torno do eixo $y$
\end{question}

%------------------------------------------------------------------------------%
\begin{resolution}{Volume rotação no eixo $y$}
  Método de cascas cilíndricas
  \[
    V = \int_a^b 2\pi \, r(x) \, h(x) \, dx
  \]


  \pause
  Raio   \tabto{20mm} \(r(x)=x\)

  \pause
  Altura \tabto{20mm} \(h(x) = f(x)\)

  \pause
  Região \tabto{20mm} \(a = 0\) até \(b = \pi\)
\end{resolution}

%------------------------------------------------------------------------------%
\begin{resolution}{Volume rotação no eixo $y$}
\begin{align*}
  \onslide<+->{V & = \int_a^b 2\pi r(x) h(x) \, dx                    }\\[2mm]
  \onslide<+->{  & = \int_{0}^{\pi} 2\pi \, x \, x \sin(x) \, dx  }\\[2mm]
  \onslide<+->{  & = 2\pi \int_{0}^{\pi} x^2 \sin(x) \, dx            }\\[2mm]
  \onslide<+->{  & = 2\pi \, F(x) \evalat{0}{\pi}                     }
\end{align*}
\end{resolution}

%------------------------------------------------------------------------------%
\begin{resolution}{Volume rotação no eixo $y$}
  \onslide<+->{\[
    F(x) = \int x^2 \sin(x) \, dx
  \]}
  \vspace{-\baselineskip}
  \begin{align*}
    \onslide<+->{u  & = x^2  & dv & =  \sin(x) dx & \hspace{3cm}~ }\\
    \onslide<+->{du & = 2xdx & v  & = -\cos(x)                    }
  \end{align*}
  \begin{align*}
    \onslide<+->{F(x)  & = -x^2 \cos(x) + 2 \int x \cos(x) \, dx }\\
    \onslide<+->{      & = -x^2 \cos(x) + 2 \, G(x)              }
  \end{align*}
\end{resolution}

%------------------------------------------------------------------------------%
\begin{resolution}{Volume rotação no eixo $y$}
  \onslide<+->{\[
    G(x) = \int x \cos(x) \, dx
  \]}
  \vspace{-\baselineskip}
  \begin{align*}
    \onslide<+->{u  & = x  & dv & = \cos(x) dx & \hspace{3cm}~ }\\
    \onslide<+->{du & = dx & v  & = \sin(x)                    }
  \end{align*}
  \begin{align*}
    \onslide<+->{G(x) & = x \sin(x) - \int \sin(x) \, dx }\\
    \onslide<+->{     & = x \sin(x) + \cos(x) + c        }
  \end{align*}
\end{resolution}

%------------------------------------------------------------------------------%
\begin{resolution}{Volume rotação no eixo $y$}
  \begin{align*}
    \onslide<+->{F(x) & = -x^2 \cos(x) + 2 \, G(x)                                }\\[2mm]
    \onslide<+->{     & = -x^2 \cos(x) + 2 \left( x \sin(x) + \cos(x) + c \right) }\\[2mm]
    \onslide<+->{     & = -x^2 \cos(x) + 2x \sin(x) + 2\cos(x) + c                }
  \end{align*}
\end{resolution}

%------------------------------------------------------------------------------%
\begin{resolution}{Solução}
  \begin{align*}
    \onslide<+->{V & = 2\pi \, F(x) \evalat{0}{\pi}                                    }\\
    \onslide<+->{  & = 2\pi \left(
                              -x^2 \cos(x) + 2x \sin(x) + 2 \cos(x)
                            \right) \evalat{0}{\pi} }\\
    \onslide<+->{  & = 2\pi \Big[
                              \big(-\pi^2 \cos(\pi) +2\pi \sin(\pi) +2\cos(\pi) \big)
                            - \big(0 + 0 + 2 \cos(0)\big)
                            \Big]                                                      }\\
    \onslide<+->{ & = 2\pi \Big[ \left( \pi^2 + 0 - 2 \right) - 2 \Big]                }\\
    \onslide<+->{ & = 2\pi \left( \pi^2 - 4 \right)                                    }\\
    \onslide<+->{ & = 2\pi^3 - 8\pi                                                    }
  \end{align*}
\end{resolution}

%------------------------------------------------------------------------------%


%------------------------------------------------------------------------------%
\section{Lista Mínima}

%------------------------------------------------------------------------------%
\begin{frame}{Lista Mínima}
  Estudar a Seção 4.1 da Apostila

  Exercícios: 1a-f, 2, 3, 9

  \vfill
  Atenção: A prova é baseada no livro, não nas apresentações
\end{frame}

%------------------------------------------------------------------------------%
\end{document}
%------------------------------------------------------------------------------%
